% --- IMPORTS ---
\documentclass[leqno]{article}
\usepackage{graphicx}
\usepackage{dsfont}
\usepackage{mathtools}
\usepackage{amssymb}
\usepackage{amsmath}
\usepackage{amsthm}
\usepackage{float}
\usepackage{ragged2e}
\usepackage{tensor}
\usepackage{fancyhdr}
\usepackage{empheq}
\usepackage[most]{tcolorbox}
\usepackage{enumitem}
\usepackage{dirtytalk}
\usepackage{marginnote}
\usepackage{microtype}
\usepackage[
  a4paper,
  left=0.8in,
  right=1in,
  top=1in,
  bottom=0.5in,
  headheight=14pt,
  headsep=0.4in
]{geometry}
\setlength{\parindent}{0cm}
% --- TikZ Packages ---
\usepackage{pgfplots}
\pgfplotsset{compat=1.18}
\usepgfplotslibrary{fillbetween, polar}
\usetikzlibrary{calc, positioning}
\usetikzlibrary{angles, quotes}
\usetikzlibrary{arrows.meta}
\usetikzlibrary{decorations.pathreplacing, decorations.markings}
\usetikzlibrary{patterns}
\usetikzlibrary{intersections}
% --- URL Packages ---
\usepackage[colorlinks=true,linkcolor=red,citecolor=magenta,urlcolor=black]{hyperref}%
\urlstyle{same}
% --- END OF IMPORTS ---

% --- CUSTOM COMMANDS ---
\RenewDocumentCommand{\marks}{o m}{%
  \IfValueTF{#1}%
    {\marginnote{\raggedright\textbf{[#2]}}[#1]}%
    {\marginnote{\raggedright\textbf{[#2]}}}%
}
\newcommand{\vspacerule}[2]{%
  \par\vspace*{#1}%
  \noindent\hrulefill\par
  \vspace*{#2}%
}
\definecolor{scarlet}{RGB}{205, 40, 75}
\newcommand{\questionitem}{%
  \item\phantomsection
  \addcontentsline{toc}{section}{Question \number\value{enumi}}%
}
\renewcommand{\contentsname}{Questions}
% --- END OF CUSTOM COMMANDS ---

% --- HEADER CONFIG ---
\pagestyle{fancy}
\fancyhead{}
\fancyfoot{}
\renewcommand{\headrulewidth}{0.9pt}
\lhead{Maclaurin Series}
\chead{}
\rhead{\href{https://danieljsmith.org}{danieljsmith.org}}
% --- END OF HEADER CONFIG ---

\begin{document}
\vspace*{20pt}
\tableofcontents
\newpage
\begin{enumerate}[
  label=\textbf{\arabic*.},
  leftmargin=*,
  topsep=0pt,
  partopsep=0pt,
  parsep=0pt
]

% ---- Question 1 ----
\questionitem

\begin{enumerate}[label=\textbf{(\alph*)}]
  \item
  \begin{enumerate}[label=(\roman*)]
    \item Given that $f(x)=e^{-2x}$, find $f'(x)$ and $f''(x)$. \marks{2}
    \item Hence, find the first three terms of the Maclaurin series for $e^{-2x}$. \marks{2} \\
  \end{enumerate}
  \item Hence, using a suitable value of $x$, show that \[e^{-1/4}\approx \frac{25}{32}\]. \marks{2} \\
\end{enumerate}

\newpage


% ---- Question 2 ----
\questionitem

\begin{enumerate}[label=(\roman*)]
  \item Use the Maclaurin series for $\sin x$ and $\cos x$ to work out the series expansion of \[\sin 2x(\cos x-\cos 3x)\] up to and including the term in $x^5$. \marks{4} \\
  \item Hence show that, in exact surd form, an approximation to the least positive root of the equation \[2\sin 2x(\cos x-\cos 3x)=x\] is \[\sqrt{\frac{4-\sqrt{10}}{12}}
\]\marks[-20pt]{3} \\
\end{enumerate}

\newpage

% ---- Question 3 ----
\questionitem

\begin{enumerate}[label=\textbf{(\alph*)}]
  \item Starting from the series given in the formulae booklet, show that the general term of the Maclaurin series for
  \[
  2(1-\cos x)-x\sin x
  \]
  is
  \[
  (-1)^r\frac{2(r-1)}{(2r)!}x^{2r}
  \]
  \marks[-20pt]{4}
  \item Show that
  \[
  \lim_{x\to 0}\left[\frac{2(1-\cos x)-x\sin x}{(1-\cos x)^2}\right]=\frac{1}{3}
  \]
  \marks[-30pt]{4}
\end{enumerate}

\newpage


% ---- Question 4 ----
\questionitem

Let $f(x)=\ln(1+\sin x)$. \\
\begin{enumerate}[label=\textbf{(\alph*)}]
  \item
  \begin{enumerate}[label=(\roman*)]
    \item Determine $f''(x)$. \marks{2}
    \item Determine the first two non-zero terms of the Maclaurin expansion for $f(x)$. \marks{3}
    \item By considering the first two non-zero terms of the Maclaurin expansion for $f(x)$, find an approximation to
    \[
    \int_0^{\frac{2}{5}} f(x)\,\mathrm{d}x
    \]
    Give your answer correct to $6$ decimal places. \marks{2} \\
  \end{enumerate}
\end{enumerate}

\newpage

% ---- Question 5 ----
\questionitem

\begin{enumerate}[label=\textbf{(\alph*)}]
  \item Given that \[y=\cosh^{-1}\left(\frac{3+x}{2+x}\right)\] show that \[(x+2)\sqrt{5+2x}\,\frac{\mathrm{d}y}{\mathrm{d}x}+1=0\] \marks{4} \\
  \item Hence find the first three terms in the Maclaurin series for $\cosh^{-1}\left(\frac{3+x}{2+x}\right)$ in the form
  \[
  a\ln\left(\frac{3+\sqrt{5}}{2}\right)+bx+cx^2
  \]
  where $a$, $b$ and $c$ are constants to be determined. \marks{5} \\
\end{enumerate}

\newpage

% ---- Question 6 ----
\questionitem

\[ y=(1+\cosh x)^n \qquad n\geq 2\] \\
\begin{enumerate}[label=\textbf{(\alph*)}]
  \item
  \begin{enumerate}[label=(\roman*)]
    \item Show that
    \[
    \frac{\mathrm{d}^2 y}{\mathrm{d}x^2}=n^2(1+\cosh x)^n-n(2n-1)(1+\cosh x)^{n-1}
    \]
    \marks[-20pt]{4}
    \item Determine an expression for
    \[
    \frac{\mathrm{d}^4 y}{\mathrm{d}x^4}
    \]
    \marks[-20pt]{2}
  \end{enumerate}
  \item Hence determine the first three non-zero terms of the Maclaurin series for $y$, giving each coefficient in simplest form. \marks{2} \\
\end{enumerate}

\newpage

% ---- Question 7 ----
\questionitem

\begin{enumerate}[label=(\roman*)]
  \item Prove that, if $y=\arctan x$, then
  \[
  \frac{\mathrm{d}y}{\mathrm{d}x}=\frac{1}{1+x^2}
  \]
  \marks[-20pt]{3}
  \item Find the Maclaurin series for $\arctan x$, up to and including the term in $x^5$. \marks{4} \\
  \item Use the result of part (ii) and the Maclaurin series for $\ln(1+t)$ to find the Maclaurin series for $\ln(1+\arctan x)$, up to and including the term in $x^4$. \marks{5} \\
\end{enumerate}

\newpage

% ---- Question 8 ----
\questionitem

\begin{enumerate}[label=\textbf{(\alph*)}]
  \item The function $f$ is defined as $f(x)=\operatorname{sech}^2 x$. \\
  \begin{enumerate}[label=(\roman*)]
    \item Show that $f^{(4)}(0)=16$. \marks{4}
    \item Hence find the first three non-zero terms of the Maclaurin series for $f(x)=\operatorname{sech}^2 x$. \marks{2} \\
  \end{enumerate}
  \item Prove that
  \[
  \lim_{x\to 0}\left(\frac{\operatorname{sech}^2 x-\cos(\sqrt{2}\,x)}{x^4}\right)=\frac{1}{2}
  \]
  \marks[-30pt]{4}
\end{enumerate}

\newpage

% ---- Question 9 ----
\questionitem

Using the Maclaurin series for $\sin x$, show that, for small values of $x$,
\[
\frac{x}{\sin x} \approx 1 + ax^2 + bx^4 + cx^6
\]
where the values of $a$, $b$ and $c$ are to be given in exact form. \marks{5}

\newpage

% ---- Question 10 ----
\questionitem

\begin{enumerate}[label=\textbf{(\alph*)}]
  \item Use the Maclaurin series expansion for $\ln(1+x)$ to show that the first three non-zero terms of the Maclaurin series expansion of
  \[
  \ln\left(\frac{1+2x}{1-x}\right)
  \]
  are
  \[
  3x-\frac{3}{2}x^2+3x^3
  \]
  \marks[-20pt]{2}
  \item Sam attempts to use the series expansion found in part (a) to find an approximation for $\ln 6$. \\

  Sam's incorrect working is shown below.
\begin{tcolorbox}[enhanced,boxrule=0.4pt,colback=white,arc=0mm,outer arc=0mm]
  \[
  \begin{aligned}
  \text{Let } \frac{1+2x}{1-x} &= 6\\
  1+2x &= 6-6x\\
  8x &= 5\\
  x &= \frac{5}{8} 
  \end{aligned}
  \] \\
  \[
  \begin{aligned}
  \text{So } \ln 6 &\approx 3\left(\frac{5}{8}\right)-\frac{3}{2}\left(\frac{5}{8}\right)^2+3\left(\frac{5}{8}\right)^3\\
  &\approx 2.02
  \end{aligned}
  \]
\end{tcolorbox}
  Explain the error in Sam's working. \marks{2} \\
  \item Use $x=\frac{1}{4}$ in the series expansion found in part (a) to find an approximation for $\ln 8$. \\

  Fully justify your answer. \marks{3} \\
\end{enumerate}

\newpage

% ---- Question 11 ----
\questionitem

\begin{enumerate}[label=\textbf{(\alph*)}]
  \item By using an appropriate Maclaurin series prove that if $0 < x < 1$ then $-\ln(1-x) > x$. \marks{2} \\
  \item Hence, by using a suitable substitution, deduce that $\ln t > 1 - \frac{1}{t}$ for $t > 1$. \marks{1} \\
  \item Using the inequality in part (b), and by making a suitable choice for $t$, determine which is greater, \[\left(\frac{5}{4}\right)^5 \quad \text{or}\quad e\] \marks{3} \\
\end{enumerate}

\newpage

% ---- Question 12 ----
\questionitem

\begin{enumerate}[label=\textbf{(\alph*)}]
  \item Write down the first three non-zero terms of the Maclaurin series for $\ln\left(\dfrac{1+2x}{1-x}\right)$. \marks{2} \\
  \item Use these three terms to show that
  \[
  \ln\left(\dfrac{10}{7}\right) \approx \dfrac{183}{512}
  \]
  \marks[-20pt]{2}
  \item Dana uses the same first three terms of the series to approximate $\ln 10$ and gets an answer of $2.67$, correct to $3$ significant figures. However, $\ln 10 = 2.30$ correct to $3$ significant figures. \\

  Explain Dana's error. \marks{2} \\
\end{enumerate}

\newpage

% ---- Question 13 ----
\questionitem

\begin{enumerate}[label=\textbf{(\alph*)}]
  \item Use the Maclaurin series expansions for $\sin x$ and $\cos x$ to determine the series expansion of
  \[
  \left(\frac{\sin(x/2)}{x/2}\right)\cos\left(\frac{x}{3}\right)
  \]
  in ascending powers of $x$, up to and including the term in $x^4$. \\

  Give each term in simplest form. \marks{3} \\
  \item Use the answer to part (a) and calculus to find an approximation, to $5$ decimal places, for
  \[
  \int_{\pi/6}^{2\pi/3} \left( \frac{1}{x}\left(\frac{\sin(x/2)}{x/2}\right)\cos\left(\frac{x}{3}\right) \right) \, \mathrm{d}x
  \]
  \marks[-30pt]{3}
  \item Use the integration function on your calculator to evaluate
  \[
  \int_{\pi/6}^{2\pi/3} \left( \frac{1}{x}\left(\frac{\sin(x/2)}{x/2}\right)\cos\left(\frac{x}{3}\right) \right) \, \mathrm{d}x
  \]

  Give your answer to $5$ decimal places. \marks{1} \\
  \item Assuming that the calculator answer in part (c) is accurate to $5$ decimal places, comment on the accuracy of the approximation found in part (b). \marks{1} \\
\end{enumerate}

\newpage

% ---- Question 14 ----
\questionitem

\begin{enumerate}[label=\textbf{(\alph*)}]
  \item Find and simplify the first five terms in the Maclaurin series for $e^x$. \marks{2} \\
  \item Hence, or otherwise, write down the first five terms in the Maclaurin series for $e^{-x}$. \marks{1} \\
  \item Use your answers, together with the identity $\cosh x + 1 = 2\cosh^2\left(\frac{x}{2}\right)$, to show that the Maclaurin series for $\cosh^2\left(\frac{x}{2}\right)$ is
  \[
  a + bx^2 + cx^4 + \cdots
  \]
  where $a$, $b$ and $c$ are rational numbers to be determined. \marks{3} \\
\end{enumerate}

\newpage

% ---- Question 15 ----
\questionitem

\begin{enumerate}[label=\textbf{(\roman*)}]
  \item Use the Maclaurin series for $\ln(1+x)$ and $\ln(1-x)$ to obtain the first three non-zero terms in the Maclaurin series for
  \[
  \ln\left(\frac{(1+x)^2}{1-2x}\right)
  \]
  State the range of validity of this series. \marks{4} \\
  \item Find the value of $x$ in the range of validity for which
  \[
  \frac{(1+x)^2}{1-2x}=2
  \]
  Hence find an approximation to $\ln 2$, giving your answer to three decimal places. \marks{4} \\
\end{enumerate}
\newpage
% ---- Question 16 ----
\questionitem

Find the Maclaurin series for \[e^{e^{x^2}-x}\] up to and including the term in $x^2$. \marks{4}

\newpage
% ---- Question 17 ----
\questionitem

\begin{enumerate}[label=\textbf{(\alph*)}]
  \item Using the definitions of $\sinh x$ and $\cosh x$, together with the Maclaurin series expansion of $e^x$, find the first three non-zero terms in the Maclaurin series expansion of $x\cosh x-\sinh x$. \marks{3} \\
  \item Hence, by replacing $x$ with $\mathrm{i}x$ in your answer to part (a), find the first three non-zero terms in the Maclaurin series expansion of $x\cos x-\sin x$. \marks{3} \\
\end{enumerate}


\end{enumerate}
\end{document}
